\section{The horizontal representation with marked bonds}\label{c7:markedgas}

The marked matrix element changes the site source and adds one
varying angle. We first establish its horizontal gas identity for
finite lists at a supported finite period. We then continue the
identity to nearby shifted lists and send the period to infinity
with those lists fixed. The estimates uniform in their sizes and
orders come in Sections~\ref{c7:shifted} and \ref{c7:mixed}; only
after those estimates do we approach the physical spread-one
lists. The auxiliary site-count limit at fixed period used for the
mixed spectral radius is stated separately below.

\subsection{The endpoint source and the finite-list regulator}

Take the combined limit of \eqref{c7:eq-L} at
$U=V=e^{\eta w_e}$. This is the only specialization of its two
endpoint variables. Write
$W_\sigma=(w_{\sigma,1},\ldots,w_{\sigma,n_\sigma})$ for the
unmarked coordinates in species $\sigma=1,2$, with
$n_\sigma\ge0$, and set $x_{\sigma,a}=e^{\eta w_{\sigma,a}}$.
The full ordered list is $(U,e^{\eta W_1},U,e^{\eta W_2})$ and
has size $M=n_1+n_2+2$. Let $P_0(W_1,W_2)$ be the all-zero term
of the finite spin sum on the unmarked lists. Empty products have
value one.

The pair factors and sign cancellation for a defect
$\delta=-s=\pm1$ at centered coordinate $v=w+i\delta/2$ are
unchanged. Before Gaussian splitting, its potentials against a
zero site in the same and opposite species are respectively
\[
 V_{d,-}(w)=q^{-1}\frac{e_1e_7}{e_0e_6},\qquad
 V_{c,-}(w)=q\frac{e_3^2}{e_4^2},\qquad
 e_k=1-q^ke^{\eta w},
\]
and $V_{\rho,+}(v-i/2)=V_{\rho,-}(v+i/2)^{-1}$.
The local factor after sign cancellation is $\sqrt c$, before
these zero-site interactions. The extra $p_0/4$ in the matched
prescription came from the opposite-species zero at the same
coordinate.

The phase $q^{-6Q}$ and the two endpoint ratios in
\eqref{c7:eq-L} contribute $\exp(i\delta D_e(v-w_e))$ per defect,
where the odd branch $D_e$, real on the real line, is determined by
\[
 e^{-iD_e(v)}
 =q^{-3}\left(\frac{e^{\eta v}-q^{11/2}}
                    {e^{\eta v}-q^{5/2}}\right)^2.
\]
It is analytic on $|\Im v|<5/2$ and tends exponentially to
$\pm\nu$ at the two real infinities, where $\nu=3\pi/4$.
Indeed
\[
 D_e'(v)=\frac{2\eta\sin(5\pi/8)}
                  {\cosh(\eta v)-\cos(5\pi/8)},\qquad
 \widehat{D_e'}(k)=4\pi\frac{\sinh(3k/2)}{\sinh4k}.
\]

Fix $0<r<1/8$. Start with distinct real unmarked coordinates
and a real endpoint coordinate, all of absolute value less than $r/8$,
and take $T>8r$. Use the Gaussian and Wick data of
Section~\ref{c7:contour}. In the current
\eqref{c7:contour-current} of species $\sigma$, put $z=1$, use the
product of $U_{\delta,T}$ translated to its own list $W_\sigma$,
and use the circle $\gamma_r$ enclosing that list. Replace
$\varphi_\sigma(v)$ by $\varphi_\sigma(v)+d_T(v)$, where
\begin{equation}\label{j:marked-kink}
 d_T(v)=\sum_{\ell\in\mathbb Z}
 \bigl[D_e(v-w_e+\ell T)-D_e(v-w_e-T/2+\ell T)\bigr]-\nu.
\end{equation}
This function has zero mean, changes sign on translation by $T/2$,
and tends to $D_e(v-w_e)$ on compact sets as $T\to\infty$.
Call the resulting Wick current $\mathcal C^L_{\sigma,T}$.
With the zero-mean source
\[
 A_\sigma(x)=\sum_{w\in W_\sigma}a_T(x-w),\qquad
 A=(A_1,A_2)^t,
\]
define the normalized regulator
\begin{equation}\label{j:marked-regulator}
 \mathfrak C_T^L=
 \frac{\mathbb E\!\left[
 e^{i\int_T(DA)\cdot\chi}
 e^{-2\pi\sum_\sigma n_\sigma\ell_\sigma/T}
 \mathcal C^L_{1,T}\mathcal C^L_{2,T}\right]}
 {\mathbb E e^{i\int_T(DA)\cdot\chi}}.
\end{equation}
The source denominator is a nonzero Gaussian exponential.

\begin{proposition}[Fixed-list marked regulator]\label{j:marked-regulated-limit}
For every $T>8r$ and fixed finite initial lists and endpoint
coordinate as above,
the Wick expansion in species $\sigma$ terminates after order
$n_\sigma$ and extends to coalescing own lists in the same
neighborhood. As $T\to\infty$ with generic initial data fixed,
\begin{equation}\label{j:marked-regulator-limit}
 \mathfrak C_T^L\longrightarrow
 \frac{N(X)^2H^{\rm flip}(X)}
 {cU^2P_0(W_1,W_2)\prod_{a\in A'\cup B'}F_0(U,x_a)^2}.
\end{equation}
This statement is initially at values where the displayed
denominator is nonzero.
\end{proposition}

\begin{proof}
The current poles inside the contour are the simple poles at its
own sites. If two residues in one species select the same site,
their first Wick fields coincide and their Pfaffian vanishes.
The kink is analytic on these contours. The same fixed-contour
argument as in Proposition~\ref{c7:regulated-pfaffian} therefore
proves termination and continuation to coalescing lists. In
particular this is a finite residue calculation at each fixed
pair of lists.

Here is the changed one-site contraction. Put
$\mathfrak a_\rho(v)=\log V_{\rho,-}(v+i/2)$, using the odd
branch zero at $v=0$. These logs have step limits. At $k>0$,
their transforms in units $2\pi/k$ are
\[
 \frac{t^{1/2}(1-t)(1+t^6)}{1-t^8},\qquad
 -\frac{2t^{7/2}(1-t)}{1-t^8},\qquad t=e^{-k},
\]
for $\rho=d,c$. Differentiate the factors and use the
hyperbolic transform from Section~\ref{c7:contour}: the exponents
in $(0,8)$ are $1.5,7.5$ versus $.5,6.5$ for the first log, and
$3.5,3.5$ versus $4.5,4.5$ for the second. Integrating the
transforms symmetrically fixes their odd primitives. The resulting
column of transforms is
$(I+gE)(\widehat a,0)^t$.

At centered argument, the own-site $U$ translate contributes
$-\delta a$ to the logarithm. The normalized Gaussian source
contributes $-\delta G_{\sigma\sigma'}D a$ per site of species
$\sigma'$. Their sum is exactly $-\delta\mathfrak a_\rho$.
The other pair and constant-mode contractions are those already
proved; in particular the difference charge is neutral, so the
singular difference-channel covariance acts only on tests
vanishing at zero frequency. Symmetric Fourier sums for the step
primitive converge after differentiation and integration from
zero. At fixed lists the winding slopes and the compensation in
\eqref{j:marked-regulator} disappear in the large-period
Gaussian sum. The local residue is $\sqrt c$; multiplying it by
$V_{c,\delta}(0)=p_0/4$ recovers the earlier matched residue.
The finite residue terms are consequently exactly those in
\eqref{c7:eq-L} after its all-zero normalization. Their convergence
is locally uniform at the fixed generic initial data, proving
\eqref{j:marked-regulator-limit}.
\end{proof}

\subsection{The changed gas data}

Write
\[
 g_+=g_d+g_c,\qquad
 l(k)=\frac{kE(k)}{1+g_+(k)E(k)},\qquad L_*:=l(0)=\frac89.
\]
Let $y_T=\partial^{-1}d_T$ have zero mean. Define $J_T^{\rm loc}$
by applying the Fourier multiplier $(l(k)-L_*)/k^2$ to $d_T'$,
using its regular value at zero. This is a periodized difference
of translates to $w_e$ and $w_e+T/2$ of a fixed exponentially
decaying function. It is analytic at position heights of absolute
value less than $5/2$ relative to either center. These facts
follow by multiplying $\widehat{D_e'}$ by the displayed symbol:
the symbol is regular at zero and meromorphic with polynomial
growth on a narrow strip about the real frequency axis.

Use the six complete particle lists \eqref{j:particle-list}, the
fused union list \eqref{c7:fused-list}, and the state gas
\eqref{c7:gas-definition}, now at the fixed post-collision
displacement $d=1/8$. The residual pair factors, local activities,
state space, and charge jumps are those of that gas at this
displacement. Its changes are the following.

\begin{enumerate}
\item In each $H$ dress of species $\sigma$, replace $B_T(v)^m$
by $\prod_{w\in W_\sigma}B_T(v-w)$. These products use only that
species' unmarked list.
\item In the integral square of \eqref{c7:state-weight}, replace
$\theta$ by $d_T(x)$; in the local Berry phase of a particle at
$x_j$, replace it by $d_T(x_j)$.
\item For an $H$ entry at height $a'$ and center $x$, with signed
power $r$, add to its dress logarithm
\begin{equation}\label{j:marked-local-dress}
 r\left[J_T^{\rm loc}(x+ia')
  -L_*\bigl(y_T(x+ia')-y_T(x)-ia'd_T(x)\bigr)\right].
\end{equation}
\end{enumerate}

The last logarithms are bounded and exponentially localized near
the two kinks, for $|\Im w_e|\le0.6$ and the listed heights. They
are evaluated on every $H$ entry, preserving the listed internal
order and the analytic fused union value. In particular the fused
particle's kink multiplier is generally not one.

For later boundary evaluation, the site part of item 1 must first
be simplified by \eqref{c7:B-recurrences}. For a particle in
species $\sigma$ it is
\begin{equation}\label{j:marked-site-factors}
\begin{array}{c|l}
 P&\displaystyle\prod_{w\in W_\sigma}B_T(x-w)\\
 M,\ F_{\rm pair}&1\\
 O_+&\displaystyle\prod_{w\in W_\sigma}B_T(x-w-i)\\
 O_-&\displaystyle\prod_{w\in W_\sigma}B_T(x-w+i)\\
 U&\displaystyle\prod_{w\in W_\sigma}B_T(x-w+i(3/2+d))^{-1}\\
 L&\displaystyle\prod_{w\in W_\sigma}B_T(x-w-i(3/2+d))^{-1}.
\end{array}
\end{equation}
The value one in the $F_{\rm pair}$ row concerns only these site
products, not item 3. Denote the resulting gas by $Z_T^L$.

\begin{proposition}[Exact marked gas]\label{j:marked-gas-identity}
Fix a sufficiently small negative pre-collision displacement
$d_-$ in the stability range, with $|d_-|<1/8$, and use
post-collision displacement $d=1/8$. For all sufficiently large
integer periods $T$, chosen independently of the finite site
counts, and for every pair of finite initial real lists and real
endpoint coordinate in the fixed neighborhood above, the gas
$Z_T^L$ converges absolutely and
\begin{equation}\label{c7:eq-M}
\begin{aligned}
 \mathfrak C_T^L
 &=c(T)Z_T^L
 \exp\left(\frac12\int_T A^tD(I+gE)A
       -i\int_T A^tD\mathbf1\,d_T+R_T^d\right),\\
 R_T^d&=-\frac{|\mathbf1|^2}{4\pi}
              \int_T d_T(l-L_*)d_T,\qquad \mathbf1=(1,1)^t.
\end{aligned}
\end{equation}
Here $c(T)$ is exactly the scalar in \eqref{c7:regulator-gas}.
The identity includes coalescing own lists by continuity. The
absolute convergence bounds may depend on the fixed lists, $T$,
and the fixed displacements. The same identity holds along $T=bn$
for a sufficiently large fixed integer block length $b$ and
positive integers $n$.
No estimate uniform as a displacement tends to zero is asserted.
\end{proposition}

\begin{proof}
We first compute each conditional current coefficient. Its loop
source is
$\mathsf A=A-i(\chi+\mathbf1d_T)$. Oscillator curvature gives
$D$ as before, and the normalized source cancels the $A,\chi$
cross term:
\[
\begin{aligned}
 &\frac12\int_T
 (A-i(\chi+\mathbf1d_T))^tD(A-i(\chi+\mathbf1d_T))
       +i\int_T A^tD\chi\\
 &\qquad=\frac12\int_T A^tDA-i\int_T A^tD\mathbf1\,d_T
        -\frac12\int_T(\chi+\mathbf1d_T)^tD(\chi+\mathbf1d_T).
\end{aligned}
\]
Completing this Gaussian shifts its mean by
$-g_+E(1+g_+E)^{-1}\mathbf1d_T$ and retains the centered
covariance $\widetilde G$ in \eqref{c7:source-removal}. Including
the source denominator, its scalar logarithm is
\[
 \frac12\int_T A^tD(I+gE)A-i\int_T A^tD\mathbf1\,d_T
 -\frac{|\mathbf1|^2}{4\pi}\int_Td_T l\,d_T.
\]
The induced common extra $H$ logarithm is
\[
 \frac{iE}{1+g_+E}d_T=-L_*y_T+J_T^{\rm loc}.
\]
The deterministic $H$ dress also has the constant
$\exp(2\pi n_\sigma/T)$. Its total exponent in species $\sigma$
is $2\pi n_\sigma\ell_\sigma/T$, because the total $H$ charge
equals $\ell_\sigma$ by \eqref{c7:gas-charge-closure}.
The winding compensation cancels it exactly.

At the post-collision endpoint, the following state rearrangement
identifies the gas defined above. The common component of $p$ is
periodic and has jumps
$-\mathbf1\cdot q_j$. Integration by parts therefore gives
\[
 \sum_j(\mathbf1\cdot q_j)y_T(x_j)
 =\int_T(\mathbf1\cdot p(x))d_T(x)\,dx.
\]
Splitting $-L_*y_T(x+ia')$ into its value at $x$, its part
linear in $a'$, and its remainder now completes the varying-angle
state square and the Berry phase. The scalar part used to complete
that square is the $L_*$ term removed from the last integral.
The remainder is exactly \eqref{j:marked-local-dress} and the
scalar left is $R_T^d$. Applying the regular multiplier
$(l-L_*)/k^2$ to $d_T'$ shows that $R_T^d$ has a finite limit as
$T\to\infty$.

We next justify summing the conditional calculation. On the first
contours the displacement is the fixed $d_-<0$, there is no
fused particle, and all $H$ heights have absolute value at most
$3/2-|d_-|$. The conditional angular trace and oscillator
calculation uses its strict strip order. Absolute Gaussian
integration uses the symmetrized pre-collision estimate of
Section~\ref{c7:stability}. At each fixed list and period, the
extra deterministic factors cost at most a constant per particle;
the quadratic count confinement sums every particle number. This
constant may depend on the finite list, whereas the supported
period threshold comes from the fixed contours and stability and
is independent of the site count.

After that integration, deform the Wick-integrated meromorphic
terms, retaining both the full heat sum and the unsplit analytic
extra $H$ logarithm $-L_*y_T(v)+J_T^{\rm loc}(v)$. The state
rearrangement above is applied at the endpoint of this move.
For the move to $d=1/8$,
the extreme $H$ height differences are at most
$3+2d=13/4<4$, the mixed $D$--external differences are at most
$5/2+d=21/8<3$, and the other outer differences stay below $2$.
The divisor checks of Section~\ref{c7:boson} therefore leave
exactly the same simple $U,L$ fusion residue. The kink dress is
analytic through it and evaluates on the fused union list.
There is no assertion about the raw Gaussian field on these
post-collision contours.

Proposition~\ref{j:stability-eighth} supplies the post-collision
absolute bound at this displacement. A bounded varying angle
changes the logarithmic modulus, relative to angle zero, by at
most linear costs in the local counts and common charge, plus
$O(T)$. The quadratic count and charge terms absorb these costs,
as well as the bounded local dress logarithms. For the initial
real lists the simplified site factors have modulus at most one;
for fixed continued lists they will instead have locally bounded
per-particle costs. The absolute estimates at the two endpoints
of the deformation and the fusion factorial identity permit
summing the finite-term identities. Fixed-count domination also
passes to coalescing own lists. This closes the proof of
\eqref{c7:eq-M}.
\end{proof}

We will also use the following precise variant. Omit $d_T$, put
$z=e^{i\theta}$ in the currents, and use the constant real angle
$\theta$ in the gas. Equation~\eqref{c7:eq-M} then holds with
the two $d_T$ terms absent, at every supported period and every
finite pair of initial real lists in $|w|<r/8$, including
coalescing lists by continuity, with bounds for fixed lists and
$\theta$ in a fixed compact set. Bounded functions of the winding
may be inserted and carried to functions of the total species
charges. In Section~\ref{c7:mixed}, the homogeneous lists
$(0,m_*)$ in this variant are taken to $m_*\to\infty$ at one
fixed supported $T$; their post-collision site factors have
modulus at most one, which supplies the needed domination in
that distinct auxiliary limit.

\subsection{Continuation and the fixed-list period limit}

At each fixed supported period, continue the finite lists through
generic coordinates whose imaginary parts, including the endpoint
coordinate, have total spread less than one. We use
$|\Im w_e|\le0.6$ and $|\Im w_{\sigma,a}|\le1$. For the final
application, take $w_e=-if$ and $w_{\sigma,a}=ih_{\sigma,a}$,
where $0\le f\le0.6$ and the $h_{\sigma,a}$ lie in the single
neighborhood of $-f,1-f$ supplied by
Proposition~\ref{j:shifted-suppression}, still with total spread
strictly less than one.

This continuation acts first on the finite residue expression
for \eqref{j:marked-regulator}. Its Gaussian source-to-defect
contractions use the rapidly converging high-frequency symbol
$\widehat G D\widehat a$; its low frequencies in the period limit
use the same odd step primitive. Defect-to-defect contractions
remain in their analytic strips, and Wick poles between generic
sites are avoided at spread less than one. The $U$ factors are
meromorphic by their defining products. Thus the finite
meromorphic identity continues from the initial real set. For
each fixed pair of counts, its finite residue expression converges
locally uniformly as $T\to\infty$ on compact subsets of the
specified generic domain of spread less than one that avoid its poles.
Indeed such a compact set has a common high-frequency decay margin
and uniform convergence of the symmetric low-frequency sums, and
there are only finitely many residue terms. Holomorphic uniqueness
therefore identifies this limit with the continued right side of
\eqref{j:marked-regulator-limit}.

On the gas side, continue the already contracted Fourier
scalars, not an uncontracted translate of $a_T$ outside its
strip. The contracted scalars need only the stated site
differences or the wider kink strip, so their period limits
converge. Use the simplified site factors
\eqref{j:marked-site-factors} before evaluating any boundary
value. They have no poles for $|\Im w_{\sigma,a}|\le1$ on the
post-collision real centers: the inverse factors may equivalently
be shifted by $3i$ using \eqref{c7:B-recurrences}. At fixed
period and fixed lists these factors are locally bounded per
particle. The stability bound therefore gives locally absolute
summability along the continuation. This proves both
\eqref{c7:eq-M} and \eqref{j:marked-regulator-limit} for the
generic continued lists just specified. It asserts no bound
uniform in the list sizes from this continuation alone.

\subsection{The growing scalar normalizer}

The exact orbital ratio in Lemma~\ref{j:shifted-nonvanishing}
now supplies the power of $M$ needed after scalar cancellation.

\begin{lemma}[Uniform shifted normalizer bound]\label{j:shifted-normalizer-bound}
For every $\varepsilon>0$, for all sufficiently large even $M$,
uniformly over real $h_j$ of spread at most one,
\begin{equation}\label{c7:eq-P}
 \frac{N(X)}{\mathcal S(X)}
 \ge \left(\frac23\right)^{M/2}M^{5/48-\varepsilon},
 \qquad x_j=e^{i\eta h_j}.
\end{equation}
\end{lemma}

\begin{proof}
Write $M=2L$ and use the law in
\eqref{j:shifted-orbital-ratio}. Under the orbital tilt put
$Q=\mathcal U D_h\mathcal U^*$ and use brackets for expectation.
If $i^-$ is the position of $-k_i$, the conjugation identity
used in Lemma~\ref{j:orbital-ward} gives
\[
 0\le\partial_{k_i}\log I_h
 =\langle Q_{ii}-Q_{i^-,i^-}\rangle
 \le\min\left\{1,\frac{k_i}{2},
 \frac{1/4}{\sum_{j\ge i}(k_i+k_j)^{-1}}\right\}.
\]
For clarity, that identity is
$\langle Q_{aa}-Q_{bb}\rangle
=((D_k)_{aa}-(D_k)_{bb})\langle|Q_{ab}|^2\rangle$.
It orders the mean diagonal entries by their test values.
The direct comparison of $i,i^-$ and $\|Q\|\le1/2$ give the
first two upper bounds. For $j\ge i$, the comparison with $j^-$
has numerator at least
$\langle Q_{ii}-Q_{i^-,i^-}\rangle$. Sum over those $j$ and
use $\sum_b\langle|Q_{ib}|^2\rangle\le1/4$ to get the last bound.

The first upper bound makes this orbital law stochastically
smaller than the pure law with weight $(k/\sinh3k)e^k$, by the
association argument in Section~\ref{c7:boundary}. Its
polynomial tail argument, applied to the first $j$ points and
then summed over $j$, shows that for
$m_0=\lfloor L^{1-\epsilon}\rfloor$, with fixed
$0<\epsilon<1$, the event
\[
 \mathcal E_L=\{k_j\le C\max(j,m_0)\text{ for every }j\}
\]
has probability $1-O(e^{-c m_0})$, uniformly in $h$. On this
event the last derivative bound is $O((\epsilon\log L)^{-1})$
for $i\le m_0$.

Fix $0<\beta<3$. Conditional on $\mathcal E_L$ and the later
points, the first $m_0$ points are stochastically smaller than
the pure law of size $m_0$ with weight
\[
 W_\beta(k)=\frac{k}{\sinh(\beta k)}\exp(\min(k^2,1))
\]
when $L$ is large. Indeed, for $k<1$ its logarithmic derivative
majorizes $(\log(k/\sinh3k))'+k/2$. For $k\ge1$ it exceeds
$(\log(k/\sinh3k))'$ by a fixed positive margin, which absorbs
the displayed $o(1)$ orbital derivative. The later
Vandermonde factors and the cutoff defining $\mathcal E_L$
decrease the relative density further. Association of the pure
law therefore proves the conditional stochastic comparison.

The function $\phi$ in \eqref{j:shifted-orbital-ratio} is
nonnegative and decreasing. Restrict to $\mathcal E_L$, discard
the later terms, and apply the conditional comparison to its
decreasing product exponential. The resulting lower bound is
$\mathbb P(\mathcal E_L)$ times the pure $W_\beta$ expectation.
Lemma~\ref{j:even-log-statistic}, rescaled by $u=\beta k/\pi$,
evaluates the latter:
\[
 \log\mathbb E_{W_\beta}\exp\left(\sum_{a=1}^{m_0}\phi(k_a)\right)
 =\frac{\beta}{\pi^2}\log m_0\int_0^\infty\phi(k)\,dk
       +o(\log m_0).
\]
To see the applicability to $W_\beta$, take the difference of that
lemma for the two even continuous tests
$\min(k^2,1)-1+\phi(k)$ and $\min(k^2,1)-1$, after rescaling.
Both decay exponentially, and the difference removes the tilt.
The integral is $5\pi^2/144$ as evaluated in
Section~\ref{c7:pfaffian}. First let $L\to\infty$ with
$\beta,\epsilon$ fixed, then let $\beta\uparrow3$ and
$\epsilon\downarrow0$. The comparisons were uniform in $h$, so
\eqref{j:shifted-orbital-ratio} gives \eqref{c7:eq-P}.
\end{proof}

\subsection{Cancellation and the interface for propagation}

We finish the scalar calculation at the fixed continued lists.
For an unmarked pair, denote its factors in
$P_0/\mathcal S(X_{\rm un})^2$ by $S_d,S_c$ for the same and
opposite species, where
$X_{\rm un}=(e^{\eta W_1},e^{\eta W_2})$. Here $\mathcal S$ is
the site prefactor of Lemma~\ref{j:shifted-nonvanishing}.
Writing $u,v$ for the site exponential variables and $C$ for its
comparison kernel at their log difference gives
\[
 S_d=\frac{(v-u)^2(v^2+u^2)}
 {(v^2-\sqrt2uv+u^2)(v^2+\sqrt2uv+u^2)C^2},\qquad
 S_c=\frac{(v-u)^2(v+u)^2}
 {(v^2+\sqrt2uv+u^2)^2C^2}.
\]
Each unmarked site versus one endpoint similarly gives, using
$F_0$,
\[
 S_e=\frac{(v-u)^2}{(v^2+\sqrt2uv+u^2)C^2}.
\]
The distinct-site quadratic contractions in \eqref{c7:eq-M}
converge to $-\log S_d$ and $-\log S_c$. Their self
exponentials are $S_d(0)^{-1/2}=2/3$ per unmarked site.
The contraction $-i\int_T A^tD\mathbf1\,d_T$ converges to
$-2\log S_e$ per unmarked site.

Here are the transform checks for these cancellations. In units
$2\pi/k$, the transforms of $\log S_\rho$ are
\[
 -\left[\frac{P_\rho}{\sinh4k}
      -\frac{\cosh3k-1}{\sinh3k}\right],
\]
where
\[
\begin{aligned}
 P_d&=\cosh4k+\cosh2k-\cosh3k-\cosh k,\\
 P_c&=\cosh4k+1-2\cosh k,\\
 P_e&=\cosh4k-\cosh k.
\end{aligned}
\]
For $d,c$ these full transforms equal
$\widehat\lambda_0(k)\widehat{\mathfrak a}_\rho(k)$.
Their negatives use $\widehat\lambda_0(-k)$ instead, exactly
as in the quadratic $A$ contraction. For $e$ the bracket is
$\tanh(3k/2)\sinh k/\sinh4k$, and
\[
 -2\widehat{\log S_e}(k)
 =\widehat\lambda_0(k)\frac{\widehat{D_e'}(k)}{k}.
\]
The kink dipole inserts the additional factor
$1-e^{-ikT/2}$; its second term averages to zero in the
periodic Riemann sum by regularity at zero and high-frequency
decay. These identities first hold for decaying logs on the real
difference line, then continue to the stated fixed-list domain.

Insert $\mathcal S(X)^2$ into
\eqref{j:marked-regulator-limit} and use \eqref{c7:eq-M}.
All distinct unmarked pairs and unmarked-to-endpoint pairs cancel
by the preceding identities. The $M-2$ self factors $2/3$
cancel $(2/3)^{-M}$ from the reciprocal square of
\eqref{c7:eq-P}, up to a constant. The reciprocal squared
Schur factor of the two endpoints, multiplied by $cU^2$, is
bounded at $U=V$ by its nonzero confluent value. The limit of
$R_T^d$ is bounded for $|\Im w_e|\le0.6$.

Consequently, fix any supported integer block length $b$. For
every fixed pair of generic imaginary lists in the common
neighborhood above, with total spread less than one including
$w_e=-if$, the limit below exists along $T=4Nb$. For every
$\varepsilon>0$ and all sufficiently large even $M$,
\begin{equation}\label{c7:eq-Z}
 |H^{\rm flip}(X)|
 \le C_\varepsilon M^{-5/24+\varepsilon}
       \left|\lim_{N\to\infty}c(4Nb)Z_{4Nb}^L\right|.
\end{equation}
The constant and threshold are uniform in $0\le f\le0.6$, the
list orders, and their finite counts. Existence of the limit
comes from the finite-list regulator limit and the convergence
of the contracted scalars, with the lists held fixed; the
uniform prefactor comes from \eqref{c7:eq-P} and the exact
cancellations. No uniform operator estimate is being inferred
from this scalar argument.

Sections~\ref{c7:shifted} and \ref{c7:mixed} now estimate the
right side uniformly for lists with the proportions in
Theorem~\ref{j:unequal}. Once that bound is available, approach
the physical spread-one heights separately for each fixed pair
of lists. Proposition~\ref{j:marked-physical} supplies the
regular limit of $H^{\rm flip}$ and its identification with
$H^{\rm poly}$ there; no period limit at the boundary is needed.
